\hspace{3mm}

\centerline{\bf \huge 6 класс}

\begin{flushright}
        \scriptsize{ $-$ Если сдаваться, не попробовав, результат не изменится.}
\end{flushright}


\problem{1} Расул Владимирович придумал систему перевода цифр в буквы и обратно, в которой каждой из десяти цифр соответствует одна уникальная буква. Расул Владимирович решил поделиться своей системой с директором Элистинского технического лицея. Директор же решила воспользоваться интересной системой и сложила в столбик два слова ЭЛМШ с ЭТЛ, получив суммарно МШМЭЛ. Определите все варианты чисел, скрывающихся за ЭЛМШ с ЭТЛ, и докажите, что других вариантов быть не может.

\problem{2}
Сугуру, Сукуна и Сатору решают олимпиаду ЭлМШ. Пока Сугуру решает $3$ задачи, Сукуна решает $4$, а пока Сукуна решает $3$ задачи, Сатору решает $5$ задач. Сколько задач решит Сугуру, если все втроём решили в общей сложности $9225$ задач?

\problem{3} 
Однажды лососи и шиншиллы выстроились в один ряд. Известно, что среди любых двадцати подряд идущих поровну лососей и шиншилл, а среди любых двадцати двух подряд идущих не поровну лососей и шиншилл. Какое максимальное количество животных могло выстроиться в один ряд?

\problem{4}
В турнире по настольному теннису участвовали $512$ человек. В каждом туре участники разбивались на пары и играли, после чего проигравшие выбывали из турнира (ничьих не бывает). После турнира, когда определился единоличный победитель, каждый из участников заявил, что в течение турнира обыграл не более одного правши. Причём все правши сказали правду, а все левши соврали. Сколько правшей могло участвовать в турнире?

\problem{5} 
Ученик за одну неделю получил 17 оценок (каждая из них – 2, 3, 4 или 5). Если сложить все оценки и разделить их на 17, то получится целое число. Докажите, что какую-то оценку он получил не более двух раз. (\textit{Агаханов Н.Х., Богданов И.И.})



